\section{Decryption Circuits}
Imagine we want to outsource checking if $a > 0$ to a third party, but we also don't want the third party to know if $a > 0$, nor do we want them to even know the value of $a$. This is a textbook use-case for HE. 

\begin{itemize}
    \item Note that $a > 0 \iff \sign(a) = 1$
    \item Derive approximation $\sign(x) \approx f(x)$ using Newton's, Householder, RK4, etc.
    \item Let's go with $\sign(x) \approx f(x) = \frac{a}{2}(3 - a^2)$
    \item Iteratively find $f_n(x) = f(f_{n-1}(x)) \to \sign(x)$ as $n \to \infty$
    \item Finally, scale results to $[0,1]$ by taking $\frac{1}{2}(\sign(x) + 1)$
    \item The reciever can decrypt the FHE result and remove noise manually i.e. determine $a > 0$ by checking if the plaintext is closer to 1 than 0.
\end{itemize}

We can therefore approximate $a > 0$ to arbitrary precision by homomorphically calculating $f_n(x)$ where $x$ is the FHE encryption of $a$ and $n$ is the number of iterations used to approximate $\sign(x)$.

% \floatname{algorithm}{Source code} % if you want to rename 'Algorithm' to 'Source code'
\begin{algorithm}[h]
  \caption{The Hello World! program in Java.}
  \label{algo:fhe:sign}
  % alternatively you may use algorithmic, or lstlisting from the listings package
  \begin{verbatim}
  
template <typename T>
Ciphertext<T> sign(CryptoContext<T> cc, Ciphertext<T> x, int iters)
{
    for(auto i = 0; i < iters; i++) {
        auto a = cc->EvalMult(x, 0.5);
        auto b = cc->EvalMult(x, x);
        b = cc->EvalSub(3, b);
        x = cc->EvalMult(a, b);
    }
    return x;
}
\end{verbatim}
\end{algorithm}

Instead of manually finding the optimal approximation of some function $f(x)$, OpenFHE provides a Chebyshev function for calculating the coefficients of a Chebyshev approximation... 

% \begin{algorithm}
% \caption{An example algorithm}
% \label{alg:example}
% \begin{algorithmic}[1] % [1] adds line numbering every line
%     \State \textbf{Input:} A set of observations $\mathcal{D}$
%     \State \textbf{Output:} A causal direction
%     \If{HSIC$_{X \to Y} <$ HSIC$_{Y \to X}$}
%         \State The causal direction is $X \to Y$
%     \ElsIf{HSIC$_{X \to Y} >$ HSIC$_{Y \to X}$}
%         \State The causal direction is $Y \to X$
%     \Else
%         \State Direction is ambiguous
%     \EndIf
%     \For{each data point in $\mathcal{D}$}
%         \State Process data
%     \EndFor
%     \State Return result
% \end{algorithmic}
% \end{algorithm}

% \begin{algorithm}[h]
%     \caption{The Hello World! program in C++.}
%     \label{algo:fhe:sign}
% \begin{lstlisting}[language=C++]
% template <typename T>
% Ciphertext<T> sign(CryptoContext<T> cc, Ciphertext<T> x, uint32_t iters = 1)
% {
%     for(uint32_t i = 0; i < iterations; i++) {
%         auto a = cc->EvalMult(x, 0.5);
%         auto b = cc->EvalMult(x, x);
%         b = cc->EvalSub(3, b);
%         x = cc->EvalMult(a, b);
%     }
%     return x;
% }
% \end{lstlisting}
% \end{algorithm}


\subsection{Sorting}
Imagine further that we have 2 numbers $a, b$ and we want to sort $\{a, b\}$ homomorphically. We can use a similar approach to before if we consider that:

\begin{equation}
    \sign(a - b) = \begin{cases} 
      \;\,1 & a > b \\
      \;\,0 & a = b \\
      -1 & a < b
   \end{cases}
\end{equation}

This allows us to sort $\{a, b\}$ by homomorphically calculating $\sign(a - b)$ and then using \emph{branchless programming} and/or boolean expressions to extract the smallest and largest values. Half the sum plus/minus half the difference will give the max/min values.

\begin{align}
    \label{eq:min}
    min &= \frac{(a + b) - |a - b|}{2} \\
    \label{eq:max}
    max &= \frac{(a + b) + |a - b|}{2}
\end{align}

This then becomes \autoref{alg:sort2}, which is shown implemented in \autoref{algo:fhe:sort_two}.

\begin{algorithm}
\caption{Homomorphic sorting of two numbers}
\label{alg:sort2}
\begin{algorithmic}[1] % [1] adds line numbering every line
    \State \textbf{Input:} Two numbers $a$ and $b$
    \State \textbf{Output:} A sorted list of the two numbers
    
    diff = a - b
    
    comp = sign(diff)
    
    abs = diff $\cdot$ comp
    
    sum = a + b

    min = (sum - abs)/2
    
    max = (sum + abs)/2
    
    \State Return $\{min, max\}$
\end{algorithmic}
\end{algorithm}


\begin{algorithm}
\caption{Sort two numbers homomorphically.}
\label{algo:fhe:sort_two}
\begin{verbatim}
template <typename T>
std::tuple<Ciphertext<T>, Ciphertext<T>> 
sort_two_unpacked(
    CryptoContext<T>    cc, 
    Ciphertext<T>       a, 
    Ciphertext<T>       b, 
    uint32_t            iterations = 2, 
    double              delta = 0.1
) {
    // Check sign(a - b)
    // diff a - b must be in [-1, 1]
    auto diff = cc->EvalSub(a, b);
    auto sign = compare(cc, diff, iterations, delta);

    // Circuit switch positions
    auto sum = cc->EvalAdd(a, b);
    auto abs = cc->EvalMult(sign, diff);

    auto min = cc->EvalMult(cc->EvalSub(sum, abs), 0.5);
    auto max = cc->EvalMult(cc->EvalAdd(sum, abs), 0.5);
    
    return std::tuple(min, max);
}
\end{verbatim}
\end{algorithm}

% \subsection{Bitonic Sorting Network}
% Sorting in this way is slow.

% Sorting networks.

% Bitonic sort can reach $\mathcal{O}(\log^2n)$, which is not quite optimal, but still. Also highly parallelizable.
